%%%%%%%%%%%%%%%%%%%%%%%%%%%%%%%%%%
% TFG — Daniel Bretaudeau — June 2026
% Built on the official UB Física template (Plantilla pel Treball Fi de Grau)
%%%%%%%%%%%%%%%%%%%%%%%%%%%%%%%%%%
\documentclass[twocolumn]{revtex4}
\usepackage{graphicx,epsfig}
\usepackage{amsmath}
\usepackage{amsfonts}
\usepackage{fancyhdr}

\newcommand{\keyword}[1]{\noindent{\bf Keywords: }{#1}\\}
\newcommand{\keywordcat}[1]{\noindent{\bf Paraules clau: }{#1}\\}
\newcommand{\sdg}[1]{\noindent{\bf SDGs: }{#1}\\}
\newcommand{\ods}[1]{\noindent{\bf ODSs: }{#1}\\}

% Kill parskip stretch + give columns a bottom sink (avoids huge gaps when
% a column is text-starved by floats)
\makeatletter
\setlength{\parskip}{0pt}
\def\@textbottom{\vskip\z@\@plus 1fil}
\makeatother

% Redefine \maketitle so that it can be used twice (for supplementary)
\makeatletter
\def\maketitle{
\@author@finish
\title@column\titleblock@produce
\suppressfloats[t]}
\makeatother

\begin{document}

\pagestyle{fancy}
\lhead{\bf Machine-learning readout of a colorimetric CO$_2$-sensing ink}
\rhead{Daniel Bretaudeau Merce}
\lfoot{Treball de Fi de Grau}
\rfoot{Barcelona, June 2026}

\title{Machine-learning readout of a colorimetric CO$_2$-sensing ink}
\author{Author: Daniel Bretaudeau Merce, dbretame17@alumnes.ub.edu}
\affiliation{Facultat de F\'{\i}sica, Universitat de Barcelona, Diagonal
645, 08028 Barcelona, Spain.}
\author{Advisor: Ismael Benito Altamirano, ismael.benito@ub.edu}

\begin{abstract}
{\bf Abstract:} A low-cost route to CO$_2$ monitoring is a colorimetric ink
whose colour tracks CO$_2$, read by miniature reflectance sensors. We characterise one ink over a 77-hour gas-chamber protocol sweeping
CO$_2$ ($\sim$100--1700 ppm), relative humidity (RH) and illumination, recording four red--green--blue--infrared pixels and NDIR
and RH references. After validating every channel against these references,
we show the blue-to-red ratio B/R responds linearly to CO$_2$
(sensitivity $\approx-5\times10^{-6}$ per ppm), with order-$10^3$ ppm
per \%RH humidity cross-sensitivity (loosely constrained) and a lag up to
tens of minutes. Humidity is recovered from the
ink alone (MAE 0.35\,\%, $R^2=0.993$). For absolute CO$_2$, referencing each colour channel to the
co-measured infrared cancels most channel-common drift: a 12-feature
linear model reaches MAE 103 ppm ($R^2=0.77$) uncorrected and, with one
offset per illumination change, 50 ppm against a 54-ppm
hold-the-calibration baseline ($R^2$ 0.71 vs 0.57) on the lamp-stable
34\,\% of test.

\keyword{colorimetric sensing, CO$_2$ sensor, optical reflectance, machine learning, regression, sensor drift}
\sdg{3, 9, 11, 13 (see page 6)}
\end{abstract}

\maketitle
\thispagestyle{fancy}

\section{Introduction}

Carbon dioxide is the standard proxy for indoor air quality and ventilation:
occupied rooms drift from the $\sim$420-ppm background to beyond
1500 ppm, levels with measurable cognitive effects \cite{satish}. The reference technology (NDIR) is accurate but bulky and
power-hungry. An attractive low-cost alternative is a \emph{colorimetric}
sensor: a dye-based ink whose colour depends on the surrounding CO$_2$
concentration, read out by nothing more than a light source and a
photodetector \cite{mills92}. The film studied here is a prototype from
a water-based, CO$_2$-selective ink family recently developed by the
host group \cite{gonzalez25}.

The transduction chain is acid--base chemistry. CO$_2$ dissolves into the
water held by the film and hydrates to carbonic acid, which dissociates and
lowers the internal pH,
\begin{equation}
\mathrm{CO_2 + H_2O \rightleftharpoons H_2CO_3 \rightleftharpoons H^+ +
HCO_3^-},
\label{eq:chem}
\end{equation}
shifting the equilibrium of a pH-indicator dye between its basic and
protonated forms. The two forms absorb in different
spectral bands, so the \emph{hue} of the reflected light (the balance
between blue and red) tracks the protonated fraction; the modulations are small ($\Delta$B/R $<1\,\%$ over the
full range), so the reflectance--absorbance relation linearises whatever
the regime, and B/R is to first order linear in that fraction. Two consequences follow. First, the
natural observable is a channel \emph{ratio}: each channel measures
illumination times film reflectance, $C_i = I_i\rho_i$, so a lamp change
that rescales every $I_i$ by a common factor cancels identically (in the
ideal case) in
$\mathrm{B/R}=\rho_{\rm B}/\rho_{\rm R}$, while any single channel is
dominated by illumination; B and R are a natural pair because, for the
indicator dyes used in CO$_2$-sensing films \cite{mills92}, the two forms
absorb at opposite ends of the visible, so $\rho_{\rm B}$ and
$\rho_{\rm R}$ shift in opposite directions. The cancellation is only as
exact as the LED's spectral stability under duty changes; Sec.~III
quantifies how well it holds in practice. Second, the protonated fraction obeys the
Henderson--Hasselbalch relation: the protonated-to-basic ratio is
proportional to $[\mathrm{H^+}]$, itself proportional to the CO$_2$
concentration in a base-buffered film \cite{mills92}, so the response is
smooth, monotonic and close to linear while that ratio remains small.
Equation~(\ref{eq:chem}) also exposes the first intrinsic weakness: water
activity enters the equilibrium directly, and the film's water content
additionally sets its swelling and the indicator's effective p$K_a$, so
humidity must interfere \cite{gonzalez25}. The second weakness is kinetic rather than
thermodynamic: reaching the equilibrium of Eq.~(\ref{eq:chem}) requires gas
dissolution and diffusion through the film (the CO$_2$ hydration step
itself adds only seconds), so the ink is \emph{slow}: effective time
constants of minutes to tens of minutes, dominated by transport.

This slowness dictated the experimental design: long quasi-static
CO$_2$ triangles ($\sim$11 h per cycle) the film can track; an RH
staircase held for hours so cross-talk maps separately; LED-duty steps to
expose what the ratios fail to cancel; and a deliberately fast final
$\sim$11 h, a stress test reserved for evaluation.

The goal is a sensor as close to \emph{self-contained} as the ink
allows: CO$_2$ and humidity from its own sixteen numbers, quantifying how
far that gets with no reference versus one maintenance offset, for a
prototype EMERGE ink.
Section III establishes that the ink senses CO$_2$ quantitatively.
Section IV develops the readout: humidity is recovered directly;
direct regression (a ladder of models from least squares to recurrent
networks under one fixed evaluation protocol) fails against the film's
slow drift, and \emph{how} each rung fails is itself informative; the
decisive step is that the sensor carries its own drift reference, the
infrared channel; one maintenance offset provides the remaining
factor of two.

\section{Experimental setup and dataset}

\subsection{Sensing element and chamber protocol}

The ink is deposited as a film and read in reflectance by four identical
OnePixel sensors \cite{benito25} (pixels 0--3), each reporting four channels, red (R),
green (G), blue (B) and infrared (IR), at 0.2 Hz under LED illumination
whose duty cycle (``lamp'') is programmable. Both multiplicities are
deliberate: separating hue from illumination needs several colour
channels (Sec.~IV shows a single one fails), and
four identical pixels give dropout redundancy, noise averaging, and a
measure of film inhomogeneity that enters the feature set below. The film sits in a sealed chamber fed by three mass-flow-controlled gas
lines (dry synthetic air, humidified air, CO$_2$) \cite{gonzalez25}, with two co-located
references: a Sensirion SCD30 NDIR \cite{scd30} for CO$_2$ and a Bosch
BME280 \cite{bme280} for RH and temperature.

\begin{figure}[t!]
\centering
\includegraphics[width=\columnwidth]{fig1_protocol.pdf}
\caption{The 77 h protocol. Top: programmed CO$_2$ setpoint (grey) and
the NDIR reference (red); the chamber tracks the setpoint ($r=0.97$) with a
systematic shortfall at the peaks, one more reason to calibrate against the
measured reference. Middle: measured relative humidity. Bottom: LED duty
cycle. Shaded bands mark the validation (grey) and test (yellow) periods;
everything before them is training data.}
\label{fig:protocol}
\end{figure}

The 77-hour protocol (Fig.~\ref{fig:protocol}): five slow triangular
CO$_2$ sweeps between $\sim$100 and $\sim$1700 ppm ($\sim$11 h each), an
RH staircase between 17 and 44 \%RH, lamp-duty steps between 20 and
80 \%, then a deliberately fast final block modulating everything
simultaneously. All streams were resampled onto a common 5-s grid
(55\,258 samples, 16 ink channels).

\subsection{Data validation}
\label{sec:validation}

Every logged channel was first cross-correlated against the two
references, and this proved essential: the acquisition log names
the gas channels in an order inconsistent with the recorded data (the
channel labelled humidity dosing tracks the NDIR CO$_2$ reference at
$r=0.97$; the one labelled CO$_2$ is the dry-air dilution line). Even the correct setpoint channel overstates the chamber (the SCD30
reaches $\sim$75\,\% of nominal, $\sim$64\,\% when the humid line runs
high): the nominal ignores part of the dilution. All results below are therefore calibrated
against the \emph{measured} references (SCD30 for CO$_2$, BME280 for RH),
and the true CO$_2$ range is the 100--1700 ppm of
Fig.~\ref{fig:protocol}.

\subsection{Learning task and splits}
\label{sec:task}

The supervised task is joint regression of the measured pair (RH, CO$_2$)
from ink reflectance only: the 16 raw channels plus 56 derived features
(per-pixel chromaticity ratios such as B/R and B/(R+G+B), and cross-pixel
means and standard deviations). Reference channels, setpoints and lamp duty
are excluded from the inputs: ratios cancel illumination by construction,
and a control adding the lamp duty as an input changes ridge-72 by
$\sim$3\,\% while degrading the IR-referenced features. Only the \emph{timing} of lamp changes is used, to
schedule that offset.

The split is \emph{chronological}: 70/15/15 train/validation/test with
4-min guard gaps (Fig.~\ref{fig:protocol}); this choice matters more than
any model choice.
Adjacent 5-s samples are nearly identical (the chamber and the film evolve
over minutes), so a \emph{random} split places near-twins of test points in
the training set and leakage masquerades as skill. Every result below is
chronological. The price is honesty about distribution shift, and here the
chronological split is also a \emph{regime} split: the protocol operates
in three modes, slow CO$_2$ triangles with an RH staircase (train), nearly
constant CO$_2$ under lamp/RH excursions (validation), and fast
simultaneous modulation with a quieter tail (test). Generalisation
therefore means transferring across operating modes as well as across
time: the test block modulates CO$_2$ on 15-min steps, faster than the
film's response (Sec.~III), its range narrows to 283--1609 ppm, and the film has aged for days. It also means validation can rank RH models but
barely CO$_2$ ones, whose variance it lacks. This mirrors a key deployment difficulty: regime shift without
retraining. As floors, predicting the training mean gives test
MAE 3.88 \%RH and 197 ppm CO$_2$; a model must beat \emph{these} rather than zero.

\section{Colorimetric characterisation}

\begin{figure}[t!]
\centering
\includegraphics[width=\columnwidth]{fig2_calibration.pdf}
\caption{Colorimetric response. (a) Pixel-averaged B/R (blue, left axis)
and reference CO$_2$ (red, right axis) during the 50\,\% lamp segment: the
ratio mirrors the full sweep; the spikes at the edges are lamp transitions.
(b) RH-corrected, mean-centred B/R versus CO$_2$ for the lamp-stable
segments with the joint fits of Eq.~(\ref{eq:linmodel}); the 50\,\% lamp
condition revisited 40 h later (red) gives a slope of the same sign and
magnitude.}
\label{fig:calibration}
\end{figure}

Within any segment of constant lamp duty (excluding the first settling hour),
the pixel-averaged blue-to-red ratio B/R follows the CO$_2$ sweeps closely
once the humidity contribution is accounted for. We fit, per segment, the
linear model
\begin{equation}
\mathrm{B/R} = \alpha\,[\mathrm{CO_2}] + \beta\,\mathrm{RH} + \gamma,
\label{eq:linmodel}
\end{equation}
and report $\alpha$ together with the partial correlation between B/R and
CO$_2$ at fixed RH (Table~\ref{tab:sens}, Fig.~\ref{fig:calibration}). Three
features stand out. First, within the scatter the response shows no
detectable curvature over the full 100--1700 ppm range, with sensitivity
$\alpha \approx -4$ to $-6 \times 10^{-6}$ ppm$^{-1}$ and partial
correlations reaching $-0.99$, consistent with the dilute limit of the
indicator equilibrium \cite{mills92}. The sign is the expected one:
more CO$_2$ acidifies the film, converting the blue-coloured basic form of
the dye (absorbing in the red) into its protonated form (absorbing in the
blue), so blue reflectance falls while red recovers and B/R drops. Second, the response \emph{persists}: the 50\,\% lamp condition
re-measured 40 h later yields a slope of the same sign and magnitude
($-6.2$ vs $-5.0\times10^{-6}$, $\sim$25\,\%). Lamp and elapsed time are
confounded, so the monotonic trend of $\alpha$ down Table~\ref{tab:sens}
need not be an illumination effect; slow gain drift of the ageing film is
the simpler reading, consistent with Sec.~IV. Third,
humidity is the dominant interferent: a global fit puts the RH term near
$10^3$ ppm of apparent CO$_2$ per \%RH, although the staircase protocol,
which holds RH nearly constant inside lamp-stable segments, leaves $\beta$
only loosely identified (its per-segment value varies in magnitude and even
sign); pinning it down needs a protocol that sweeps RH and CO$_2$ jointly.
Any quantitative use of the ink thus requires simultaneous knowledge of RH,
which, as shown next, the ink itself provides.

\begin{table}[b!]
\centering
\caption{Per-segment joint fits of Eq.~(\ref{eq:linmodel}) in lamp-stable
windows (first settling hour excluded). $r_p$ is the partial correlation of
B/R with CO$_2$ at fixed RH. The initial 70\,\% lamp sweep (1--13 h) is reported but excluded from
the summary: it is the only segment at that lamp level and the only one
with a fresh, still-equilibrating film, so the two effects are confounded
($\alpha=-7.5\times10^{-6}$, $r_p=-0.68$).}
\label{tab:sens}
\vspace*{1mm}
\begin{tabular}{ccccc}
\hline\hline
Lamp (\%) & $t$ (h) & CO$_2$ (ppm) & $\alpha$ ($10^{-6}$/ppm) & $r_p$ \\
\hline
60 & 14--24.5 & 108--1680 & $-4.3$ & $-0.77$ \\
50 & 25.5--36 & 102--1668 & $-5.0$ & $-0.95$ \\
40 & 37--46.5 & 99--1486  & $-5.7$ & $-0.99$ \\
30 & 47.5--55 & 93--1070  & $-5.9$ & $-0.91$ \\
50 & 75--76.7 & 282--556  & $-6.2$ & $-0.83$ \\
\hline\hline
\end{tabular}
\end{table}

The dynamics are not instantaneous: cross-correlating RH-corrected B/R
with the NDIR gives a $\sim$3-min apparent lag in slow sweeps but up to
$\sim$44 min in the fast block. The film acts as a low-pass filter whose apparent lag grows with
excitation speed; fluctuations faster than it are attenuated, not
tracked.

\section{Predicting RH and CO$_2$ from reflectance}

\subsection{Models and evaluation}

Table~\ref{tab:ladder} reports every model on identical inputs,
splits and metrics, with train next to test error so the
\emph{generalisation gap} is visible per family;
Fig.~\ref{fig:allmodels} (appendix) shows every fit over the full run.
Instantaneous-model hyperparameters were tuned on validation only; it
carries almost no CO$_2$ variance (Sec.~II.C), so the LSTM configuration
is a representative architecture.

\begin{table}[t!]
\centering
\caption{The model ladder, curated to one representative per lesson
(Fig.~\ref{fig:allmodels} in the appendix shows every fit). Train/test MAE
for both targets (RH in \%, CO$_2$ in ppm) and test Pearson $r$ for CO$_2$;
$R^2_{\rm te}<0$ for all CO$_2$ entries except the IR-referenced ridge
($R^2=0.77$) and the bottom block
($R^2$ = 0.57/0.65/0.71/0.68). $N_p$ = trainable parameters
(``---'' = nonparametric). Floors (predict the train mean): 3.88 \%RH,
197 ppm. Bottom block: one 30-min offset per lamp
change, evaluated on lamp-stable test windows (34\,\% of test);
LSTM $r$/offset-MAE: five-seed means (other cells one
archived run); MLP: three-seed means.}
\label{tab:ladder}
\vspace*{1mm}
\resizebox{\columnwidth}{!}{%
\begin{tabular}{lrcccccl}
\hline\hline
 & & \multicolumn{2}{c}{RH MAE} & \multicolumn{2}{c}{CO$_2$ MAE} & & \\
\cline{3-4}\cline{5-6}
Model & $N_p$ & tr & te & tr & te & $r_{\rm te}$ & What it teaches \\
\hline
Ridge, 72 features   & 146  & 0.17 & 0.88 & 119 & 403  & 0.76 & physics-guided features help \\
Ridge + high-pass 4 h& 578  & 0.04 & 0.77 & 44  & 297  & 0.43 & drift lives at low frequency \\
Ridge, 12 (R,G,B)/IR & 26   & 0.27 & 1.74 & 179 & {\bf 103} & {\bf 0.90} & IR reference cancels drift \\
MLP 128$\times$64, reg. & 18k & 0.23 & 1.34 & 26  & 521  & 0.70 & regularisation alone did not fix drift \\
LSTM, 6-h coarse     & 21k  & 0.14 & 0.69 & 19  & 654  & 0.81 & long context: shape, not scale \\
\hline
Hold calibration value & 0   & ---  & ---  & --- & 54   & 0.79 & the bar recalibration sets \\
B/R + RH + recal.    & 3    & ---  & ---  & --- & 54   & 0.82 & ties the bar, tracks better \\
(R,G,B)/IR + offset  & 26   & ---  & ---  & --- & {\bf 50} & {\bf 0.91} & the deployable readout \\
LSTM 6-h + offset    & 21k  & ---  & ---  & --- & 48 & 0.82 & regime-sensitive experiment \\
\hline\hline
\end{tabular}}
\end{table}

\subsection{Humidity}

\begin{figure}[t!]
\centering
\includegraphics[width=\columnwidth]{fig3_predictions.pdf}
\caption{Test-block predictions. (a) RH from ink reflectance alone.
(b) The 6-h-context LSTM tracks the CO$_2$ oscillations ($r=0.81\pm0.03$
over seeds; best of five shown, drawn after its 6-h warm-up) but loses the
absolute scale; for display its output is affinely rescaled, which is
precisely the calibration a deployment must supply. (c) The same LSTM
with one 30-min offset per lamp change (red, drawn where evaluated)
recovers absolute CO$_2$ in the lamp-stable windows. Fluctuations faster
than the film's $\sim$44-min apparent lag are smoothed by the ink
itself.}
\label{fig:predictions}
\end{figure}


Against the 3.88 \%RH floor, an RH-tuned ridge (8-h high-pass window,
selected per target on validation) reaches {\bf MAE 0.35 \%RH,
$R^2=0.993$} on the test block (Fig.~\ref{fig:predictions}a). Physically
expected: water activity enters Eq.~(\ref{eq:chem}) directly.
An ablation pins the mechanism: red alone fails below the floor
(4.4 \%RH; one band is absorbance times illumination), one pixel's four
colours recover 0.60 \%RH, and sixteen raw channels do \emph{worse} than
four (1.04). The RH signal is spectral, and the ink itself measures its dominant
interferent.

\subsection{CO$_2$ by direct regression}

The same recipes collapse on CO$_2$: every plain learner ends below the
197-ppm floor, killed by the slow baseline drift of the ageing film.
The failures are systematic and each teaches something
(Table~\ref{tab:ladder}). Least squares on the 16 raw channels simply
memorises the training period: 87 ppm on train, 1549 on test. Engineered
ratios repair part of that (403 ppm) and causal high-pass features a bit
more (297 ppm at a 4-h window), but the CO$_2$ signal rides a baseline
whose drift overlaps the protocol cycles, so no static feature set
removes it. More capacity makes things worse: growing an unregularised
MLP from 4.8k to 2.2M parameters drives training error down while test
error climbs from 671 to 1030 ppm, and proper regularisation does not
rescue it ($521\pm67$ ppm, $r=0.70$, three seeds). Sequence models come
closest: an LSTM reading six hours of 60-s-averaged context tracks the
shape of the oscillations better than anything else ($r=0.81\pm0.03$
over five seeds, Fig.~\ref{fig:predictions}b) yet still posts MAE
$\approx$650 ppm, because no context length we could test (up to 8 h)
supplies the absolute anchor. The pattern repeats in every family,
linear to recurrent: none infers the drift from the visible channels;
added capacity only memorises it.

\subsection{CO$_2$ with an infrared reference}

The way out is not a bigger model but a better ratio. The dye barely
touches the \emph{infrared} channel, so IR shares the slow ageing of
emitter, detector and film with the colour channels while staying blind
to the gas, and the division cancels what they share: from train to test, B/IR shifts by only 0.26
train-$\sigma$, against 1.09 for B/R and 0.69 for raw B. A ridge on the
twelve (R,G,B)/IR ratios reaches {\bf MAE 103 ppm, $R^2=0.77$, $r=0.90$
with no correction at all}, the only model that beats the constant
floor unaided. Its test error is even \emph{lower} than its training
error (103 vs 179 ppm); no miracle, the test block just spans a
narrower range, and relative to each block's spread both errors are the
same $\approx$0.4$\sigma$. What the ratio does not cancel are the lamp
steps (IR barely follows the duty cycle), and the remaining error is
made of exactly that. The fitted weights mostly sit on the red-channel
ratios with opposing signs across pixels, which points to between-pixel
common-mode rejection (collinear coefficients are not individually
meaningful), and tiny MLPs only lose: the relation is essentially
linear.

\subsection{CO$_2$ with one offset per illumination change}

After an illumination change no ink signal supplies the absolute
anchor; the standard remedy is one offset, measured against a reference in
the first 30 min of each lamp segment (samples excluded; evaluated
windows = 34\,\% of test, 283--623 ppm). The relevant baseline is strict: holding the calibration value
already gives 54 ppm ($R^2=0.57$); the corrected ridge (122) and MLP
(79, one seed) fall short, and the three-parameter
Eq.~(\ref{eq:linmodel}) ties it with better tracking ($R^2=0.65$), with
measured or ink-derived RH alike.
The IR-referenced ridge beats the bar: {\bf 50 ppm, $R^2=0.71$,
$r=0.91$} (Fig.~\ref{fig:predictions}c), twelve features,
selected by validation (Fig.~\ref{fig:summary}). The six-hour LSTM matches it numerically ($47.7\pm1.2$ ppm over five
seeds) but only while training matches the test regime; in the single
refit run of Sec.~IV.F it collapses, and it is reported as an \emph{experiment} on temporal
context. Without the offset, Eq.~(\ref{eq:linmodel})
fails like everything else (171 ppm): the residual drift could only be
removed by measuring it.

Three caveats bound these numbers. The offsets are computed against the
co-located NDIR reference, so a deployment needs a reference channel or a
reference-free scheme (e.g.\ baseline autocalibration \cite{vergara}).
All results come from one film in one 77-h run. The single-ratio B/R noise after 60-s averaging,
$\sim$$1.1\times10^{-4}$ ($\approx$22 ppm at the measured sensitivity),
is an indicative resolution scale; in practice drift dominates the
error budget. Chamber temperature (27.3--31.2 $^\circ$C) correlates
weakly with the residuals ($r=0.26$), so it is not the dominant term.

These limits also say what the ink \emph{is} for. As a ppm meter it
needs slow environments and an occasional anchor; as an \emph{indicator},
neither: with no correction at all, the IR-referenced readout reaches
AUC 0.98 at 800 ppm and 0.99 at 1000 ppm (accuracies 92\,\% and 96\,\%
against majority-class baselines of 79\,\% and 88\,\%). That is the
indoor-ventilation use the host group's monitoring platform targets
\cite{gonzalez25,benito25}: a low-cost element asking only whether the
level is above a threshold.

\subsection{Adding the validation block to training}

Refitting on train+validation after selection (injecting lamp/RH
variation at constant CO$_2$) improves every instantaneous model
markedly (ridge+HP 297$\to$137 ppm, MLP 521$\to$155, IR ridge
103$\to$86 without correction and 50$\to$45 ppm with offset, in this single
configuration),
while the recurrent model collapses in our single refit run
($r$ 0.81$\pm$0.03 $\to$ 0.05, one seed), consistent with the added
block teaching a wrong temporal regime. This suggests the binding constraint is the data itself: one film,
one run, five similar sweeps, one confounder block.

\section{Conclusions}
\begin{itemize}
\item The ink is a genuine CO$_2$ sensor: in lamp-stable conditions
B/R shows no detectable deviation from linearity over 100--1700 ppm
($\alpha \approx -5\times10^{-6}$ ppm$^{-1}$, partial correlations
to $-0.99$); the slope keeps sign and magnitude after 40 h
($\sim$25\,\% gain drift, absorbed by the IR reference).

\item Humidity is both the dominant interferent (order $10^3$ ppm per
\%RH, loosely identified) and accurately recoverable from the ink
itself on this film (MAE 0.35 \%RH, $R^2=0.993$).

\item Absolute CO$_2$ readout is limited by slow drift rather than
model capacity; most of that drift is channel-common and cancelled by the
infrared reference: a 12-feature (R,G,B)/IR linear model reaches 103 ppm
($R^2=0.77$) without correction and 50 ppm with one 30-min offset per
illumination change (45 in a single refit on all pre-test data); all
results are one-film, reference-anchored.

\item A central limitation is the \emph{speed} of the CO$_2$ changes:
the apparent lag grows from $\sim$3 to $\sim$44 min under fast
modulation. Fluctuations faster than this are physically unrecoverable, a
floor on test error distinct from model failure (Appendix); this sets the
application envelope (slowly varying air); next: drift-stable references,
an orthogonal RH$\times$CO$_2$ sweep, a multi-film repeat. In threshold mode the absolute-anchor requirement disappears
(AUC 0.98 at 800 ppm with no correction): the ventilation-indicator regime.

\item Methodologically, validating every channel against independent
references proved decisive: it exposed a channel-labelling inconsistency
and redefined the calibration targets.
\end{itemize}

\begin{acknowledgments}
I thank my advisor, Ismael Benito Altamirano, and Cristian F\`abrega
with the EMERGE collaboration for ink, chamber and measurements.
\end{acknowledgments}

\vspace*{-7mm}
\begin{thebibliography}{99}

\bibitem{satish} U. Satish {\it et al.}, ``Is CO$_2$ an indoor pollutant?
Direct effects on human decision-making performance'', Environ. Health
Perspect. {\bf 120}, 1671 (2012).

\bibitem{mills92} A. Mills, Q. Chang and N. McMurray, ``Equilibrium studies
on colorimetric plastic film sensors for carbon dioxide'', Anal. Chem.
{\bf 64}, 1383 (1992).

\bibitem{gonzalez25} M. Gonz\'alez-G\'omez, I. Benito-Altamirano,
J. D. Prades, O. Casals and C. F\`abrega, ``Ultralow-cost and selective
water-based colorimetric ink for indoor CO$_2$ monitoring'', IEEE Sens.
J. {\bf 25}, 18722--18730 (2025).

\bibitem{benito25} I. Benito-Altamirano {\it et al.}, ``When colorimetry
meets the cloud: IoT-enabled health monitoring at home via
highly-specific gas sensing'', Proc. AINA-2025, Springer (2025).

\bibitem{scd30} Sensirion AG, {\it SCD30 CO$_2$ sensor datasheet}
(2020).

\bibitem{bme280} Bosch Sensortec, {\it BME280 sensor datasheet} (2022).

\bibitem{vergara} A. Vergara {\it et al.}, ``Chemical gas sensor
drift compensation using classifier ensembles'', Sens. Actuators B
{\bf 166}, 320 (2012).

\end{thebibliography}

\clearpage

\title{Lectura per aprenentatge autom\`atic d'una tinta colorim\`etrica
sensible al CO$_2$}

\begin{abstract}
{\bf Resum:} Una via de baix cost per monitorar el CO$_2$ \'es una tinta
colorim\`etrica el color de la qual canvia amb la concentraci\'o de CO$_2$,
llegida per sensors de reflect\`ancia miniaturitzats. Caracteritzem una
tinta d'aquest tipus dins d'una cambra de mescla de gasos durant un protocol
de 77 hores que escombra el CO$_2$ ($\sim$100--1700 ppm), la humitat relativa
(HR) i la il$\cdot$luminaci\'o, enregistrant quatre p\'ixels de reflect\`ancia
vermell--verd--blau--infraroig juntament amb una refer\`encia NDIR
(CO$_2$) i una d'HR. Despr\'es de validar el conjunt de dades amb les refer\`encies
mesurades, mostrem que la ra\'o blau/vermell respon linealment al CO$_2$ amb
sensibilitat $\approx -5\times10^{-6}$\,ppm$^{-1}$, una sensibilitat creuada
a l'HR de l'ordre de $10^3$ ppm per \%HR, i un retard de resposta de minuts a
desenes de minuts. Comparem models de calibratge de capacitat creixent, des dels m\'inims
quadrats fins a xarxes neuronals recurrents. La humitat
es recupera nom\'es amb la tinta (MAE 0,35\,\%, $R^2=0{,}993$); per al CO$_2$
absolut, referenciar cada canal de color a l'infraroig co-mesurat
cancel$\cdot$la la major part de la deriva \`optica lenta: un model lineal
de 12 caracter\'istiques assoleix MAE 103 ppm ($R^2=0{,}77$) sense cap
correcci\'o, i 50 ppm amb un ajust puntual per canvi
d'il$\cdot$luminaci\'o, sobre el 34\,\% del bloc de test estable en
il$\cdot$luminaci\'o.

\keywordcat{detecci\'o colorim\`etrica, sensor de CO$_2$, reflect\`ancia
\`optica, aprenentatge autom\`atic, regressi\'o, deriva de sensors}
\ods{3, 9, 11, 13}
\end{abstract}

\maketitle
\thispagestyle{fancy}

\onecolumngrid
\vspace{-0.5cm}
\section*{Objectius de Desenvolupament Sostenible (ODSs o SDGs)}
\vspace{-0.2cm}
\begin{center}
\begin{tabular}{|l|c|l|c|}
\hline
1. Fi de la pobresa & & 10. Reducci\'o de les desigualtats &\\ \hline
2. Fam zero	               &  &	11. Ciutats i comunitats sostenibles	& X \\  \hline
3. Salut i benestar	      & X &	12. Consum i producci\'o responsables& \\  \hline
4. Educaci\'o de qualitat   &  &	13. Acci\'o clim\`atica	& X \\  \hline
5. Igualtat de g\`enere	      &  &	14. Vida submarina	& \\  \hline
6. Aigua neta i sanejament &  &	15. Vida terrestre	& \\  \hline
7. Energia neta i sostenible &  &	16. Pau, just\'icia i institucions s\`olides	& \\  \hline
8. Treball digne i creixement econ\`omic &  & 17. Alian\c{c}a pels objectius	& \\  \hline
9. Ind\'ustria, innovaci\'o, infraestructures & X &  & \\
\hline
\end{tabular}
\end{center}
\vspace{0.1cm}
El contingut d'aquest TFG es relaciona amb l'ODS 3 (Salut i benestar), fita
3.9, ja que el monitoratge de la qualitat de l'aire interior mitjan\c{c}ant
sensors de CO$_2$ de baix cost contribueix a reduir l'exposici\'o de les
persones a ambients mal ventilats. Es relaciona tamb\'e amb l'ODS 9
(Ind\'ustria, innovaci\'o, infraestructures), fita 9.5, perqu\`e desenvolupa i
caracteritza una tecnologia de sensors innovadora i de baix consum; amb
l'ODS 11 (Ciutats i comunitats sostenibles), fita 11.6, perqu\`e un
monitoratge ambiental assequible i ubic \'es una eina per a edificis i
ciutats m\'es saludables; i amb l'ODS 13 (Acci\'o clim\`atica), fita 13.3, ja
que millora la capacitat de mesura del CO$_2$, el principal gas amb efecte
d'hivernacle d'origen antropog\`enic.

\clearpage
\twocolumngrid
\appendix
\section{SUPPLEMENTARY MATERIAL}

\begin{figure}[b!]
\centering
\includegraphics[width=\columnwidth]{fig6_summary.pdf}
\caption{Summary: test CO$_2$ MAE of the main models (log scale).
Direct regression never beats the constant-prediction floor; the
IR-referenced ridge beats it with no correction (103 ppm); with one 30-min
offset per lamp change, it and the six-hour LSTM tie at 50/48 ppm.}
\label{fig:summary}
\end{figure}

\begin{figure*}[t!]
\centering
\includegraphics[width=0.96\textwidth]{fig5_allmodels.pdf}
\caption{Four models over the full 77 h run, every one fitted on the
training block only: prediction (red, clipped to the axis range) against
the NDIR reference (grey); validation (grey band) and test (yellow band)
marked. Each model fits the training period;
the bands reveal who generalises. The large MLP and Ridge-72 lose the
scale to drift; the LSTM keeps the shape; only the uncorrected
IR-referenced ridge stays on the measured trace throughout. Quoted MAEs
are the test values of Table~\ref{tab:ladder}.}
\label{fig:allmodels}
\end{figure*}

\begin{table}[b!]
\centering
\caption{Configurations behind Table~\ref{tab:ladder}. All models:
standardised inputs, chronological splits; ridge/MLP hyperparameters
chosen on validation.}
\label{tab:hyper}
\vspace*{1mm}
\resizebox{\columnwidth}{!}{%
\begin{tabular}{lll}
\hline\hline
Model & Inputs & Key settings \\
\hline
Ridge & 72 feat. & $\alpha\in\{10,10^2,10^3\}$ on val \\
Ridge + HP & 72 + 216 high-pass & HP windows 1/2/4/8 h, + 5/15-min smooth \\
MLP & 72 feat. & ReLU 128$\times$64, drop 0.2, wd $10^{-4}$, early stop \\
LSTM (coarse) & 16 raw, 60-s MA, 1/min & $W{=}360$ (6 h), hidden 64, 40 epochs \\
B/R + RH & pixel-mean B/R, RH & least squares on train; 30-min offset \\
\hline\hline
\end{tabular}}
\end{table}


{\bf Training details} (Table~\ref{tab:hyper}){\bf .} The ridge penalty ($\in\{10,10^2,10^3\}$) and the
high-pass window were selected on the validation block only. The LSTMs use a
single recurrent layer of 64 units trained with Adam (learning rate
$10^{-3}$, 25--40 epochs, batch 512) on standardised inputs and targets;
training windows are strided by 30 s, test windows by 5 s. The coarse-context
variant first applies a 60-s moving average to the 16 channels, resamples at
one point per minute, and reads windows of 360 such points (6 h); a 12-h
window is not measurable because the test block is only 11.5 h long.
Window-normalised LSTMs (each input window z-scored independently, 30/60-min
windows) end at CO$_2$ MAE 386/299 ppm with $r=0.12/0.10$. All models receive identical
chronological splits with 4-min guard gaps; stochastic models (MLPs,
LSTMs) use a fixed random seed and each configuration is a single run,
except the coarse-context LSTM and its recalibrated variant (five
seeds: shape $r$ 0.76--0.85; corrected MAE 46.7--49.9 ppm) and the
regularised MLP (three seeds). All rolling windows are trailing
(causal).

{\bf Recalibration policy.} The 30-min window after each lamp change is used
to fit a single additive offset per model and segment; those samples are
excluded from all reported metrics, and windows shorter than 30 min
contribute no evaluated samples. Evaluated coverage is 34\,\% of the test
period (3.9 h), spanning 283--623 ppm. Reducing the calibration window to
10 min degrades the physical model to MAE $\sim$100 ppm; extending it to
60 min shrinks the evaluated windows and also degrades the result.

{\bf What else was tried.} (0) Raw-window LSTMs (5--30 min of 5-s
history) reach only 574--582 ppm: the relevant history exceeds any
affordable raw window, which motivated the coarse-context design.
(i) A derivative-predicting RNN integrated
from each anchor (dead reckoning, 10-seed ensemble) won validation
decisively but lost on test (61 ppm): a derivative bias of only
0.025 ppm\,s$^{-1}$ is invisible in dynamics-rich validation yet fatal in
flat test windows: regime mismatch can silently invert rankings.
(ii) Kinetic features ($d/dt$, $d^2/dt^2$ of ratios and channels at 1--60
min scales) all lost to their static controls, and lead-compensating the
film's lag was rejected at every $\tau$ from 3 to 90 min: differentiation
amplifies the RH signal; lag compensation is not the lever. (iii) A
PatchTST-style transformer (105k params, val-selected, 3 seeds) lost to the
LSTM on every metric (85 vs 48 ppm corrected; shape $r\approx-0.4$).
(iv) Adversarial RH-decoupling via gradient reversal only hides RH from
the online head: a fresh probe still recovers RH at $R^2\approx0.93$, and
forcing real invariance collapses CO$_2$ in lockstep: the two signals
share the protonation pathway and cannot be separated in this ink.

{\bf Train+validation refit.} Full numbers of the Sec.~IV.F refit:
ridge-72 403$\to$250 ($R^2$ still $<0$); ridge+HP 297$\to$137
($R^2=0.48$); MLP 521$\to$155 ($R^2=0.33$); IR ridge 103$\to$86
($R^2=0.80$, $r=0.92$) and, with the offset, 50$\to$45 ppm
($R^2=0.76$); B/R+RH+offset 54$\to$53; LSTM 6-h shape $r$ 0.81$\to$0.05
and 48$\to$93 ppm with offset. The validation block teaches confounder
robustness to instantaneous models and a wrong temporal regime to
recurrent ones.

{\bf Reproducibility.} The full pipeline (feature construction, every
model, figure and verification scripts) is a documented set of Python
scripts, available on request; total training time
is under five minutes on CPU for all linear models plus
$\sim$10 min on one GPU for the LSTM variants.

\end{document}
